% Template LaTeX - Pôster IX Simpósio de Iniciação Científica / UNIFEI
% Dimensões e tipografia ajustadas ao modelo PPTX oficial.
% Versão compatível com pdfLaTeX.
% Compilar com:
%   pdflatex banner_simposio_pdflatex.tex
%
% O modelo PPTX utiliza Verdana nos títulos e Arial no corpo.
% Para garantir portabilidade com pdfLaTeX, este template utiliza
% Bera Sans como família sans-serif substituta.
%
\documentclass[portrait]{a0poster}
%
\usepackage[T1]{fontenc}
\usepackage[utf8]{inputenc}
\usepackage[portuguese]{babel}
% Bera Sans family supplied by the TeX distribution.
\usepackage[scaled]{berasans}
\renewcommand{\familydefault}{\sfdefault}
\usepackage[paperwidth=90cm, paperheight=120cm, left=6.7cm, right=6.7cm, top=17.2cm, bottom=13.0cm]{geometry}
%
\usepackage{indentfirst}
\usepackage{graphicx}
\usepackage{eso-pic}
\usepackage{xcolor}
\usepackage{ragged2e}
\usepackage{multicol}
\usepackage{titlesec}
\usepackage{microtype}
%
% ------------------------------------------------------------
% FONTES -- versão compatível com pdfLaTeX
% ------------------------------------------------------------
% O PPTX exige Arial no corpo e Verdana nos títulos.
% pdfLaTeX não acessa diretamente fontes TrueType/OpenType do sistema.
% Bera Sans (fvs) é utilizada tanto no corpo quanto nos títulos,
% evitando dependência de fontes instaladas no sistema operacional.
\newcommand{\posterbody}{\fontfamily{fvs}\selectfont}
\newcommand{\posterheading}{\fontfamily{fvs}\selectfont}

% ------------------------------------------------------------
% CABEÇALHO E RODAPÉ
% Mantidos como no pré-modelo fornecido.
% image2.png = cabeçalho; image3.png = rodapé.
% ------------------------------------------------------------
\AddToShipoutPictureBG{%
  \AtPageUpperLeft{%
    \raisebox{-12.25cm}{\includegraphics[width=\paperwidth]{image2.png}}%
  }%
}
\AddToShipoutPictureBG{%
  \AtPageLowerLeft{%
    \raisebox{-0.15cm}{\includegraphics[width=\paperwidth]{image3.png}}%
  }%
}

% ------------------------------------------------------------
% PARÂMETROS GERAIS DO PÔSTER
% ------------------------------------------------------------
\pagestyle{empty}
\setlength{\parskip}{0.55em}
\setlength{\emergencystretch}{2em}

% No PPTX, as duas colunas têm aproximadamente 36 cm de largura,
% separadas por cerca de 4 cm.
\setlength{\columnsep}{4.2cm}
\setlength{\columnseprule}{0pt}

% Corpo: Arial 28 pt, parágrafo simples, justificado.
\newcommand{\postertext}{%
  \posterbody
  \fontsize{29}{42}\selectfont
  \justifying
  \setlength{\parindent}{2.0cm}%
}

% Títulos das seções: Verdana 44 pt, negrito.
%
\newcommand{\sectiontitle}[1]{%
  \scalebox{1.08}[1.04]{%
    \posterheading
    \fontsize{42}{52}\selectfont
    \bfseries
    #1%
  }%
}

\titleformat{\section}
  {\normalfont}
  {}{0pt}{\sectiontitle}
% ------------------------------------------------------------
% COMANDOS PARA TÍTULO/AUTORES
% ------------------------------------------------------------
\newcommand{\postertitle}[1]{%
  \begin{center}
    \parbox{55cm}{%
      \centering
      \posterheading
      \fontsize{72}{80}\selectfont
      \bfseries
      \MakeUppercase{#1}\par
    }%
  \end{center}
}

\newcommand{\posterauthors}[1]{%
  \begin{center}
    {\posterheading\fontsize{34}{38}\selectfont\bfseries\itshape #1\par}
  \end{center}
}

\newcommand{\posterinstitution}[1]{%
  \begin{center}
    {\posterheading\fontsize{36}{38}\selectfont #1\par}
  \end{center}
}

\begin{document}

% ------------------------------------------------------------
% BLOCO SUPERIOR
% ------------------------------------------------------------
% O PPTX inicia o título aproximadamente a 18,1 cm do topo.
\vspace*{-1.75cm}

\postertitle{O título do trabalho deve ser em fonte Verdana,
tamanho 66, letras maiúsculas, em negrito}

\vspace{1.55cm}

\posterauthors{SOBRENOME, Autor 1., SOBRENOME, Autor 2\ldots  Verdana, 32, negrito, centralizado}

\vspace{-0.95cm}

\posterinstitution{INSTITUIÇÃO DE ORIGEM}

\vspace{1.90cm}

% ------------------------------------------------------------
% DUAS COLUNAS
% ------------------------------------------------------------
\begin{multicols}{2}
\postertext

\section*{INTRODUÇÃO}

Este modelo de pôster descreve o estilo a ser usado no VIII Simpósio de
Iniciação Científica da Unifei.

A fonte a ser utilizada deve ser Arial. Aconselhamos fonte com tamanho
mínimo 28. Pode-se deixar um recuo em cada parágrafo, a critério do autor.

O texto deve ser digitado em parágrafo simples, justificado.

O título deve estar em caixa alta, centralizado, fonte Verdana, tamanho 66
em negrito. O espaço disponível para o título deverá ser mantido.

A caixa de texto do título do trabalho, bem como dos autores devem ser
colocadas de forma que fiquem centralizadas, no meio do Pôster.

O título e o nome dos autores deverão ser idênticos ao do resumo submetido.

\vspace{2.0cm}
\section*{METODOLOGIA}

O nome do(s) autor(es) deve(m) estar centralizado(s), com fonte Verdana,
tamanho 32, negrito e itálico. Para mais de um autor, os nomes deverão estar
dispostos na mesma linha e separados por vírgulas. Inicialmente, apresenta-se
o sobrenome, em caixa alta, seguido do restante do nome.

\vspace{1.2cm}
\section*{RESULTADOS OBTIDOS E ANÁLISE}

As dimensões padrão do Pôster devem ser de 90 cm de largura, com 120 cm de
altura.

Não será permitida apresentação por terceiros (não autores). Os trabalhos
de Iniciação Científica devem ser apresentados pelo discente, acompanhado
pelo orientador.

O autor do trabalho deverá permanecer junto ao pôster durante o período da
apresentação, junto com o respectivo orientador, para responder às questões
dos interessados.

\columnbreak

Continuação do campo referente aos resultados obtidos e análises.

Poderão ser inseridos gráficos, tabelas e figuras.

\vspace{10.0cm}
\section*{CONCLUSÕES}

Incluir as conclusões do trabalho. Manter fonte Arial, tamanho mínimo 28.

\vspace{10.0cm}
\section*{AGRADECIMENTOS E FINANCIAMENTOS}

No caso de estudantes da Unifei, é obrigatório constar no pôster
agradecimentos à universidade e, no caso de alunos bolsistas, ao órgão
financiador da bolsa. Outros agradecimentos são opcionais.

\vspace{2.0cm}
\section*{REFERÊNCIAS}

As referências devem ser feitas respeitando-se as normas definidas pela
ABNT. Somente deverão ser apresentadas as referências citadas no referido
pôster.

\end{multicols}

\end{document}
